\documentclass[journal]{IEEEtran}

% ---- Packages ----
\usepackage{amsmath}
\usepackage{amssymb}
\usepackage{booktabs}
\usepackage{graphicx}
\usepackage{algorithm}
\usepackage{algpseudocode}
\usepackage{subcaption}
\usepackage{multirow}
\usepackage{xcolor}
\usepackage{cite}
\usepackage{url}
\usepackage{hyperref}   % Must be loaded LAST among these packages

\graphicspath{{figures/}}

% ---- Title ----
\title{Adaptive Spectral Defense for Real-Time Recovery of\\
  Deep-Learning-Based Modulation Classifiers\\
  Under Adversarial Attack}

% ---- Authors ----
\author{%
  Author~One,~\IEEEmembership{Member,~IEEE,}
  Author~Two,~\IEEEmembership{Member,~IEEE,}
  and~Author~Three%
  \thanks{Manuscript received \today.}%
}

\markboth{IEEE Trans.\ on Cognitive Commun.\ and Netw.}%
{Author \MakeLowercase{\textit{et al.}}: Adaptive Spectral Defense}

\begin{document}

\maketitle

% ---- Abstract ----
\begin{abstract}
Deep-learning-based automatic modulation classifiers (AMC) achieve high
accuracy under benign channel conditions but collapse under adversarial
perturbations crafted by optimization-based attacks such as Carlini--Wagner
(CW) and Elastic-net Attack (EAD).  Existing defenses require model
retraining or provide insufficient accuracy recovery in the signal domain.
We propose Adaptive-K v2, a unified frequency-domain defense that filters
adversarial perturbations without retraining the classifier.  The method
combines three mechanisms in a single shared-FFT pass: spectral flatness
routing for wideband modulations, SNR-adaptive K-cap to balance filtering
aggressiveness with signal preservation, and per-sample magnitude-knee
estimation to adapt to each modulation's spectral width.
We evaluate against five attacks (CW, EAD-L1, EAD-EN, FGSM, PGD) on
RML2016.10a using the pretrained Adaptive Wavelet Network (AWN) classifier
and compare against five calibrated classical filters (Kalman, Wiener,
Savitzky-Golay, Gaussian, FIR) and Randomized Smoothing.
Optimization attacks are devastating: CW reduces accuracy to 35.4\% and
EAD-L1 to 6.6\%.  Adaptive-K v2 recovers CW accuracy to 71.3\%
(+35.9~pp) and EAD-L1 to 70.8\% (+64.2~pp), outperforming the best
classical filter (Kalman) by 6.2~pp on CW and 8.5~pp on EAD-L1.
On gradient attacks, Adaptive-K v2 recovers PGD to 43.4\% (+10.9~pp)
and FGSM to 52.0\% (+3.9~pp).  Classical filters are counterproductive
on FGSM, reducing accuracy below the undefended baseline despite
per-SNR calibration.  These results demonstrate that a single adaptive
spectral method outperforms all classical signal processing baselines
across both optimization and gradient attack families.
The defense adds less than 0.1~ms latency overhead per batch and degrades
clean accuracy by $<$0.5~pp from the 92.6\% baseline.
\end{abstract}

% ---- Keywords ----
\begin{IEEEkeywords}
adversarial machine learning, automatic modulation classification,
input transformation defense, spectral filtering, RF security
\end{IEEEkeywords}

% ---- Sections ----
\section{Introduction}
\label{sec:intro}

% TODO: content in Plan 02
% Key points to cover:
% - AMC importance in cognitive radio and spectrum monitoring
% - DNN-based classifiers (AWN) achieve state-of-the-art accuracy
% - Adversarial attacks (CW, EAD, FGSM, PGD) severely degrade performance
% - Existing defenses (adversarial training, feature squeezing) are either slow or require retraining
% - Our contribution: detect-recover-classify pipeline with FFT Top-K recovery
% - Summary of results and paper organization

\section{Related Work}
\label{sec:related}

% TODO: content in Plan 02

\subsection{Deep Learning for Automatic Modulation Classification}
% TODO: AMC classifiers — O'Shea CNN, CLDNN, ResNet-AMC, AWN, LSTM-based

\subsection{Adversarial Attacks on AMC Systems}
% TODO: gradient-based attacks on AMC — Lin threats, Flowers SigGuard, Silvaco survey, stealthy attacks

\subsection{Input-Transformation Defenses}
% TODO: Guo input transformations, FFT-based defenses, randomized smoothing (Cohen), feature squeezing (Xu)

\subsection{Classical Signal Processing Filters as Defenses}
% TODO: Kalman, Wiener, Savitzky-Golay, Gaussian, FIR in the context of adversarial robustness

\section{System Model and Threat Model}
\label{sec:system}

\subsection{Signal Representation and Dataset}

We consider complex baseband IQ signals represented as real-valued tensors
$\mathbf{x} \in \mathbb{R}^{2 \times T}$, where the first dimension indexes
the in-phase (I) and quadrature (Q) channels and $T$ is the number of time
samples. Each channel captures a different aspect of the baseband signal; their
joint representation encodes the amplitude and phase of the modulated carrier.
Following the RML2016.10a dataset convention~\cite{otoole2016rml}, we use
$T = 128$ samples, corresponding to 16 symbols at 8 samples-per-symbol with
a root-raised-cosine pulse shape.

The RML2016.10a dataset contains 220,000 IQ bursts drawn from 11 modulation
classes: BPSK, QPSK, 8PSK, QAM16, QAM64, PAM4, AM-DSB, AM-SSB, WBFM, GFSK,
and CPFSK. Each burst is labeled with the modulation type and the signal-to-noise
ratio (SNR) at which it was recorded, spanning $-20$ to $+18$~dB in 2~dB steps.
At high SNR ($\geq 0$~dB), the AWN classifier~\cite{li2023awn} achieves above
91\% overall accuracy across all 11 classes.

\subsection{AWN Classifier Architecture}

The Adaptive Wavelet Network (AWN)~\cite{li2023awn} takes batched IQ tensors of
shape $[N, 2, T]$ and predicts one of $C = 11$ modulation classes. The forward
pass has three stages:

\textbf{Convolutional feature extraction.} A 2D convolution with kernel size
$[2, 7]$ integrates the I and Q channels, yielding a feature map of shape
$[N, 64, 1, T]$, which is then squeezed to $[N, 64, T]$. A second 1D
convolution with batch normalization~\cite{ioffe2015batch} and LeakyReLU
extracts temporal features, producing $[N, 64, T]$ feature maps.

\textbf{Adaptive wavelet decomposition.} The core module applies a
learnable lifting scheme~\cite{li2023awn} that recursively splits the
feature map into even and odd subsequences, applies learnable Predictor ($P$)
and Updator ($U$) operators to produce approximation coefficients $\mathbf{c}$
and detail coefficients $\mathbf{d}$, and repeats for \texttt{num\_levels}
levels. The detail regularization loss is:
\begin{equation}
  \mathcal{R}_d = \sum_{\ell=1}^{L} \lambda_d \,\|\mathbf{d}_\ell\|_2^2,
  \quad
  \mathcal{R}_c = \lambda_c \,|\bar{\mathbf{c}}|,
\end{equation}
where $\bar{\mathbf{c}}$ is the mean of the final approximation. The total
training loss is:
\begin{equation}
  \mathcal{L} = \mathcal{L}_{\mathrm{CE}} + \mathcal{R}_d + \mathcal{R}_c,
\end{equation}
where $\mathcal{L}_{\mathrm{CE}}$ is the cross-entropy classification loss.

\textbf{Attention and classification.} Squeeze-excitation (SE)
attention~\cite{hu2018squeeze} computes per-level weights for the multi-scale
approximation and detail coefficient tensors. The attention-weighted features
are pooled and passed through two fully-connected layers with dropout and
LeakyReLU activations to produce the final class logits. The model is trained
using the Adam optimizer~\cite{kingma2015adam} with early stopping.

\subsection{Attack Models}

We evaluate four attacks that span the gradient-based and optimization-based
attack families. All attacks use \texttt{torchattacks}~\cite{kim2020torchattacks}
with per-sample minmax normalization: each IQ sample is mapped to $[0, 1]$ using
$\tilde{x} = (x - x_{\min}) / (x_{\max} - x_{\min})$ before attack generation
and mapped back afterward. This makes epsilon values interpretable relative to
the signal range.

\textbf{FGSM.} The Fast Gradient Sign Method~\cite{goodfellow2015fgsm}
computes a single gradient step:
\begin{equation}
  \mathbf{x}_{\mathrm{adv}} = \mathbf{x} + \varepsilon \cdot
    \mathrm{sign}\bigl(\nabla_{\mathbf{x}}\,\mathcal{L}_{\mathrm{CE}}
    (\mathbf{x}, y)\bigr),
\end{equation}
where $\varepsilon$ bounds the $L_\infty$ perturbation magnitude. FGSM
is computationally efficient but produces suboptimal adversarial examples.

\textbf{PGD.} Projected Gradient Descent~\cite{madry2018pgd} iterates
multiple gradient steps with projection onto the $\varepsilon$-ball:
\begin{equation}
  \mathbf{x}^{(t+1)} = \Pi_{\mathcal{B}_\varepsilon(\mathbf{x})}
    \!\left(\mathbf{x}^{(t)} + \alpha\,\mathrm{sign}\bigl(
    \nabla_{\mathbf{x}}\,\mathcal{L}_{\mathrm{CE}}(\mathbf{x}^{(t)}, y)
    \bigr)\right),
\end{equation}
where $\Pi$ denotes projection and $\alpha$ is the step size. PGD with a
random start is regarded as a strong first-order attack~\cite{croce2020reliable}.

\textbf{CW ($L_2$).} The Carlini--Wagner attack~\cite{carlini2017towards}
solves the constrained optimization:
\begin{equation}
  \min_{\boldsymbol{\delta}}\; \|\boldsymbol{\delta}\|_2 + c\, f(\mathbf{x} + \boldsymbol{\delta}),
\end{equation}
where $f(\cdot)$ is a surrogate misclassification objective and $c$ balances
perturbation norm against attack success. CW distributes adversarial energy
broadly across spectral frequencies, making it resistant to simple band-limited
filtering. We use $c = 10$, 200 optimization steps.

\textbf{EAD ($L_1$ and EN).} The Elastic-Net Attack on Deep Neural
Networks~\cite{chen2020ead} adds an $L_1$ regularization term to the CW
objective:
\begin{equation}
  \min_{\boldsymbol{\delta}}\; \beta\|\boldsymbol{\delta}\|_1 +
    \|\boldsymbol{\delta}\|_2 + c\, f(\mathbf{x} + \boldsymbol{\delta}),
\end{equation}
where $\beta$ controls the sparsity of the perturbation. We evaluate both the
$L_1$ variant (EADL1, $\beta$ large) and the elastic-net variant (EADEN,
$\beta$ moderate). EADL1 produces sparse perturbations concentrated in few
spectral components, while EADEN produces a mixture of sparse and distributed
perturbations.

\subsection{Threat Model}

\textbf{Attacker capabilities.} We consider a white-box adversary with full
knowledge of the AWN classifier architecture, trained weights, and input
normalization procedure. The attacker can query the classifier and compute
exact gradients. This is the strongest adversary assumption and establishes
a worst-case defense baseline.

\textbf{Attacker goal.} The adversary seeks untargeted misclassification: any
prediction other than the true modulation class is considered a success. The
attacker does not attempt to cause the classifier to predict any specific wrong
class.

\textbf{Attacker constraint.} Perturbations are bounded by the chosen
attack norm: $\|\boldsymbol{\delta}\|_\infty \leq \varepsilon$ for FGSM/PGD,
$\|\boldsymbol{\delta}\|_2 \leq$ (implicitly, via CW minimization) for CW, and
elastic-net bounded for EAD. We evaluate $\varepsilon \in \{0.05, 0.1\}$ in
the per-sample minmax-normalized space (equivalent to 5--10\% of the
per-sample signal range).

\textbf{Defense-unaware model.} We assume the attacker does not adapt to the
defense mechanism: adversarial examples are generated against the undefended
AWN classifier. This models the realistic deployment scenario where the defense
is applied transparently at the receiver and is not publicly disclosed. Adaptive
attacks against the defense are left for future work.

\textbf{Defender capabilities.} The defender can preprocess each IQ burst
before classification, with access only to the received signal. The defender
does not know the true modulation class, the attack type, or the attack
parameters. The defense must operate in real time: latency must be sub-millisecond
per sample to meet the requirements of practical AMC deployment contexts,
including ITU/FCC spectrum monitoring stations, CBRS Environmental Sensing
Capability (ESC) nodes, and military electronic support measures (ESM)
systems~\cite{proakis2008digital}. In each of these contexts, AMC drives
downstream decisions (spectrum allocation, threat identification) that require
near-real-time throughput, typically at 25--100~Msps IQ digitization rates with
burst processing latency budgets of 1--10~ms per channel.

\section{Proposed Defense Methods}
\label{sec:method}

\subsection{Control-Plane Attack Insight}

Before describing our defenses, we motivate the frequency-domain approach by
establishing the nature of adversarial attacks on AMC. We conduct a two-track
experiment to separate \emph{classification failure} from \emph{waveform
corruption}. In Track~A, we apply CW and EAD attacks to real RML2016.10a
signals and record classification accuracy. In Track~B, we generate synthetic
bursts with known payload bits and CRC-8 checksums using a full TX/RX chain
(pulse shaping, AWGN channel, matched filtering, pilot-based equalization),
transfer the adversarial perturbation vectors from Track~A, and evaluate CRC
pass rate under four demodulation conditions:

\begin{enumerate}
  \item Clean+Oracle: clean signal with true modulation known;
  \item Clean+AMC: clean signal, modulation predicted by AWN;
  \item Adv+Oracle: adversarial signal with true modulation known;
  \item Adv+AMC: adversarial signal, modulation predicted by AWN.
\end{enumerate}

Table~\ref{tab:crc} shows CRC pass rates at SNR~18~dB under CW attack for
eight digital modulations. The Adv+Oracle column reveals that the underlying
signal remains largely intact under adversarial perturbation: CRC pass rates
stay above 71\% when the correct demodulator is applied. The catastrophic
failure manifests only in the Adv+AMC column (0--39\%), where the classifier
selects the wrong demodulator.

\begin{table}[t]
\caption{CRC Pass Rate (\%) at SNR 18~dB --- CW Attack. Adversarial
  perturbation drops AMC accuracy to near 0\% (Adv+AMC) but CRC integrity
  under oracle demodulation (Adv+Oracle) remains above 71\%, confirming
  that attacks target the classification decision, not the physical signal.}
\label{tab:crc}
\centering
\small
\begin{tabular}{lcccc}
\toprule
Mod & Clean+Orc & Clean+AMC & Adv+Orc & Adv+AMC \\
\midrule
BPSK  & 100.0 & 97.0 & 100.0 & 39.0 \\
QPSK  & 100.0 & 98.5 & 100.0 &  1.5 \\
8PSK  & 100.0 & 98.0 & 100.0 &  0.5 \\
QAM16 & 100.0 & 94.5 &  99.5 &  0.5 \\
QAM64 &  96.0 & 94.0 &  76.0 &  0.5 \\
PAM4  & 100.0 & 98.5 &  71.0 &  2.0 \\
CPFSK & 100.0 & 100.0 & 100.0 & 18.5 \\
GFSK  &  99.5 & 99.0 &  99.5 &  0.0 \\
\bottomrule
\end{tabular}
\end{table}

This \emph{control-plane attack} insight motivates our frequency-domain
approach: adversarial perturbations are spectral artifacts that can in principle
be attenuated by frequency-domain filtering without destroying the underlying
modulation waveform.

\subsection{FFT Top-\texorpdfstring{$K$}{K} Filtering}

The fundamental building block of our defense is FFT Top-$K$ filtering,
which retains the $K$ largest-magnitude spectral components and zeros all
others. Formally, for a single IQ sample $\mathbf{x} \in \mathbb{R}^{2 \times T}$:

\begin{enumerate}
  \item \textbf{Normalize:} $\tilde{\mathbf{x}} = (\mathbf{x} + x_{\mathrm{off}})
    / x_{\mathrm{scale}}$, mapping amplitudes to $[0, 1]$.
  \item \textbf{FFT:} $\mathbf{X} = \text{rFFT}(\tilde{\mathbf{x}})$, producing
    a complex spectrum $\mathbf{X} \in \mathbb{C}^{2 \times (T/2+1)}$.
  \item \textbf{Top-$K$ mask:} For each channel $c \in \{I, Q\}$, compute
    magnitudes $|\mathbf{X}_c|$ and construct binary mask $\mathbf{m}_c$
    that retains the $K$ largest entries and zeros the rest.
  \item \textbf{Inverse FFT:} $\hat{\tilde{\mathbf{x}}} =
    \text{irFFT}(\mathbf{X} \odot \mathbf{m})$.
  \item \textbf{Denormalize:} $\hat{\mathbf{x}} = \hat{\tilde{\mathbf{x}}}
    \cdot x_{\mathrm{scale}} - x_{\mathrm{off}}$.
\end{enumerate}

In compact notation:
\begin{equation}
  \hat{\mathbf{x}} = \text{irFFT}\!\bigl(\mathrm{TopK}\bigl(\text{rFFT}(\tilde{\mathbf{x}}),\, K\bigr)\bigr),
  \label{eq:topk}
\end{equation}
where the normalization and denormalization are implicit. The rationale is that
adversarial perturbations, which are added in the amplitude domain, tend to
distribute energy across the full spectrum; retaining only the dominant $K$
components removes much of this diffuse adversarial energy while preserving the
modulation-carrying low-frequency components.

\subsection{Spectral-Gated Defense}
\label{sec:sg}

FFT Top-$K$ filtering is effective for narrowband modulations (BPSK, QPSK,
QAM, etc.) but completely fails for wideband modulations such as AM-SSB,
whose spectral energy is spread uniformly across frequencies. Top-$K$ filtering
of AM-SSB destroys the signal itself, reducing recovery accuracy to 0\% even
when $K$ is generous.

We exploit a robust spectral feature, \emph{spectral flatness}, to detect
wideband signals without a separate classifier. Spectral flatness is defined as
the ratio of the geometric mean to the arithmetic mean of the power spectrum:
\begin{equation}
  \text{SF}(\mathbf{x}) = \frac{\exp\!\bigl(\tfrac{1}{N}\sum_{k}\ln |X_k|^2\bigr)}
                               {\tfrac{1}{N}\sum_{k}|X_k|^2},
  \label{eq:flatness}
\end{equation}
where $X_k$ are the rFFT coefficients and $N = T/2 + 1$ is the number of
positive-frequency bins. For narrowband modulations at SNR~18~dB, spectral
flatness is below 0.04 (all energy concentrated in few bins). For AM-SSB,
flatness is approximately 0.55 (energy spread uniformly across frequencies).
This gap of more than $18\times$ provides a robust routing threshold with
negligible misclassification between the two routes.

Algorithm~\ref{alg:gated} presents the spectral-gated defense. The key
efficiency insight is that the FFT computation is \emph{shared}: the same
FFT coefficients are used for both the flatness check and the Top-$K$
operation, requiring no additional computation.

\begin{algorithm}[t]
\caption{Spectral-Gated Defense}
\label{alg:gated}
\begin{algorithmic}[1]
\Require Input $\mathbf{x} \in \mathbb{R}^{N \times 2 \times T}$,
  Top-$K$ parameter $K$, quantization levels $L$, flatness threshold $\tau$
\Ensure Defended signal $\hat{\mathbf{x}}$
\State Normalize: $\tilde{\mathbf{x}} \leftarrow (\mathbf{x} + x_{\mathrm{off}}) / x_{\mathrm{scale}}$
\State $\mathbf{X} \leftarrow \text{rFFT}(\tilde{\mathbf{x}})$
  \hfill \textit{// shared FFT for all samples}
\For{each sample $i = 1, \ldots, N$}
  \State $\text{SF}_i \leftarrow \text{SpectralFlatness}(\mathbf{X}_i)$
    \hfill \textit{// Eq.~\eqref{eq:flatness}}
  \If{$\text{SF}_i > \tau$}
    \hfill \textit{// wideband (AM-SSB)}
    \State $\hat{\tilde{\mathbf{x}}}_i \leftarrow \text{Quantize}(\tilde{\mathbf{x}}_i, L)$
  \Else
    \hfill \textit{// narrowband (all others)}
    \State $\hat{\tilde{\mathbf{x}}}_i \leftarrow \text{irFFT}(\mathrm{TopK}(\mathbf{X}_i, K))$
      \hfill \textit{// reuse FFT}
  \EndIf
\EndFor
\State Denormalize: $\hat{\mathbf{x}} \leftarrow \hat{\tilde{\mathbf{x}}} \cdot x_{\mathrm{scale}} - x_{\mathrm{off}}$
\State \Return $\hat{\mathbf{x}}$
\end{algorithmic}
\end{algorithm}

Default parameters are $K = 20$, $L = 32$, and $\tau = 0.4$.
The quantization routing recovers AM-SSB from 0\% to 20--65\% accuracy under
CW and EAD attacks, a modulation where Top-$K$ alone achieves 0\% recovery.

\subsection{Adaptive-\texorpdfstring{$K$}{K} Defense}
\label{sec:adaptivek}

Fixed-$K$ filtering applies the same spectral truncation to all signals
regardless of their inherent bandwidth, which can over-filter wideband
digital modulations or under-filter narrowband signals under low-SNR attacks.
The Adaptive-$K$ defense selects the number of retained spectral components
per sample based on the spectral magnitude distribution.

Algorithm~\ref{alg:adaptivek} describes the per-sample $K$ selection procedure.
The key idea is that the spectral magnitude profile of a clean modulated signal
has a characteristic shape: a few dominant components carry most of the energy,
followed by a rapid decay toward noise floor. The \emph{knee point} of this
profile---where the cumulative energy first exceeds a threshold $\eta$---defines
the natural $K$ for that signal. Adversarial perturbations, by distributing
energy across many bins, push the knee point to larger $K$ values; the adaptive
method corrects for this by anchoring $K$ to the cumulative energy threshold
rather than a fixed count.

\begin{algorithm}[t]
\caption{Adaptive-$K$ Spectral Defense}
\label{alg:adaptivek}
\begin{algorithmic}[1]
\Require Input $\mathbf{x} \in \mathbb{R}^{N \times 2 \times T}$,
  energy threshold $\eta \in (0, 1)$, $K_{\min}$, $K_{\max}$
\Ensure Defended signal $\hat{\mathbf{x}}$
\State Normalize: $\tilde{\mathbf{x}} \leftarrow (\mathbf{x} + x_{\mathrm{off}}) / x_{\mathrm{scale}}$
\State $\mathbf{X} \leftarrow \text{rFFT}(\tilde{\mathbf{x}})$
\For{each sample $i = 1, \ldots, N$}
  \State $\mathbf{m}_i \leftarrow |\mathbf{X}_i|$ \hfill \textit{// per-channel magnitudes, flattened}
  \State Sort $\mathbf{m}_i$ descending: $\mathbf{m}_i^{\downarrow}$
  \State $E_{\mathrm{total}} \leftarrow \sum_k (\mathbf{m}_i^{\downarrow}[k])^2$
  \State $K_i \leftarrow \min\{k : \sum_{j=1}^{k}(\mathbf{m}_i^{\downarrow}[j])^2
    \geq \eta \cdot E_{\mathrm{total}}\}$
  \State $K_i \leftarrow \mathrm{clip}(K_i,\; K_{\min},\; K_{\max})$
  \State $\hat{\tilde{\mathbf{x}}}_i \leftarrow \text{irFFT}(\mathrm{TopK}(\mathbf{X}_i, K_i))$
\EndFor
\State Denormalize: $\hat{\mathbf{x}} \leftarrow \hat{\tilde{\mathbf{x}}} \cdot x_{\mathrm{scale}} - x_{\mathrm{off}}$
\State \Return $\hat{\mathbf{x}}$
\end{algorithmic}
\end{algorithm}

Default parameters are $\eta = 0.95$, $K_{\min} = 10$, $K_{\max} = 40$.
For narrowband modulations (BPSK) at high SNR, the selected $K$ is typically
10--15; for wideband digital modulations (QAM64, GFSK), it is 20--35.
On the RML2016.10a dataset, adaptive-$K$ achieves the highest overall
accuracy among all evaluated defenses, outperforming fixed Top-$K$ by 1--3\%
on CW attacks and 3--5\% on EAD attacks across SNR 0--18~dB.

\subsection{Classical Filter Baselines}
\label{sec:classical}

To contextualize the performance of our frequency-domain adaptive methods,
we evaluate five classical signal processing filters as adversarial defense
baselines. These represent the standard signal processing toolkit that a
communications engineer would apply before considering ML-specific defenses.
All filters are calibrated per-SNR via grid search over their hyperparameters
on a held-out validation subset of RML2016.10a, maximizing post-filter
classification accuracy on the undefended model.

\textbf{Kalman filter.} The Kalman filter~\cite{kalman1960new} models each IQ
channel as a first-order random walk corrupted by observation noise, with
process noise covariance $Q$ and observation noise covariance $R$ as tunable
parameters. The scalar filter recursively estimates the signal state using
prediction and update steps. For AMC defense, $R$ is calibrated to the noise
floor of adversarial perturbations. Kalman filtering is optimal for linear
Gaussian noise but was not designed for structured adversarial perturbations;
our experiments test whether its denoising properties generalize.

\textbf{Wiener filter.} The Wiener filter~\cite{wiener1949extrapolation}
computes the MMSE estimate in the frequency domain by applying a
signal-to-noise weighting $H_k = S_{ss,k} / (S_{ss,k} + S_{nn,k})$ to each
spectral bin, where $S_{ss,k}$ and $S_{nn,k}$ are the signal and noise power
spectral densities at frequency $k$, estimated from calibration data.
The Wiener filter assumes stationary signal and noise statistics; adversarial
perturbations violate the stationarity assumption.

\textbf{Savitzky-Golay filter.} The Savitzky-Golay filter~\cite{savitzky1964smoothing}
smooths each IQ channel by fitting a polynomial of degree $d$ to a sliding
window of $w$ samples using local least squares. The filter preserves the
shape of peaks and pulses better than standard lowpass filters, making it
suitable for preserving QAM constellation transitions. Parameters $d$ and
$w$ are calibrated via grid search.

\textbf{Gaussian filter.} A Gaussian kernel $h[n] = \exp(-n^2 / 2\sigma^2)$
is convolved with each IQ channel, implementing a smooth lowpass response with
cutoff determined by the standard deviation $\sigma$. The filter is implemented
as a GPU-native \texttt{F.conv1d} operation for efficient batch processing,
avoiding CPU roundtrip latency.

\textbf{FIR lowpass filter.} A finite impulse response (FIR) lowpass filter
with normalized cutoff frequency $f_c$ and filter order $M$ is applied to
each IQ channel~\cite{proakis2006digital}. The filter coefficients are computed
via windowed-sinc design (Hamming window). Like the Gaussian filter, the FIR
filter is implemented as a depthwise \texttt{F.conv1d} convolution for GPU
efficiency.

We also evaluate \textbf{randomized smoothing}~\cite{cohen2019certified} as a
certified robustness baseline. Randomized smoothing adds isotropic Gaussian
noise $\mathcal{N}(0, \sigma^2 I)$ to the input and takes the majority vote
over $k$ noisy copies. The certified radius is $\sigma \Phi^{-1}(p_A)$, where
$p_A$ is the probability of the top class under smoothing. We use $\sigma = 0.1$
and $k = 100$ copies, which is the lowest latency configuration that provides
non-trivial certificates. All five classical filters and randomized smoothing
are calibrated and evaluated under identical experimental conditions.

\subsection{Computational Complexity}
\label{sec:complexity}

The computational cost of each defense is dominated by its core operation.
FFT Top-$K$ (and adaptive-$K$) requires $\mathcal{O}(T \log T)$ for the real
FFT, $\mathcal{O}(T)$ for magnitude computation, $\mathcal{O}(T \log K)$ for
top-$K$ selection via partial sort, and $\mathcal{O}(T \log T)$ for the inverse
FFT. The spectral flatness computation adds $\mathcal{O}(T)$. Total per-sample
cost is $\mathcal{O}(T \log T)$, dominated by the FFT.

Among the classical filters, Gaussian and FIR convolutions cost
$\mathcal{O}(T \cdot M)$ per channel, where $M$ is the filter length (typically
17--33 taps); both are implemented as GPU-native depthwise convolutions with
negligible latency. Savitzky-Golay uses a precomputed coefficient vector and has
the same $\mathcal{O}(T \cdot w)$ complexity as a FIR filter. Kalman filtering
requires a sequential state update loop, with $\mathcal{O}(T)$ operations per
sample but poor parallelism (the state update at time $t$ depends on time $t-1$),
necessitating a CPU fallback and incurring higher per-batch latency despite lower
asymptotic complexity. Wiener filtering applies per-bin frequency-domain
weighting in $\mathcal{O}(T \log T + T)$ using a precomputed gain vector.
Randomized smoothing requires $k$ forward passes through the classifier ($k=100$
by default), making it $100\times$ more expensive than a single inference.
Detailed latency measurements and throughput comparisons are reported in
Section~\ref{sec:results}.

\section{Experimental Setup}
\label{sec:setup}

\subsection{Dataset and Model}

We evaluate on the RadioML 2016.10a dataset~\cite{otoole2016rml}, a widely
used benchmark for automatic modulation classification comprising 220{,}000 IQ
signal frames spanning 11 modulation classes: eight digital (BPSK, QPSK, 8PSK,
QAM16, QAM64, PAM4, CPFSK, GFSK) and three analog (AM-DSB, AM-SSB, WBFM).
Each frame consists of 128 complex-valued samples represented as a $2\times128$
real tensor of in-phase (I) and quadrature (Q) components.  The dataset covers
SNR values from $-20$\,dB to $+18$\,dB in 2\,dB increments, yielding 20 SNR
levels and 1{,}000 frames per (modulation, SNR) cell.

We partition the data using stratified sampling that preserves the joint
(modulation, SNR) distribution across splits: 70\% training, 15\%
validation, and 15\% test.  All reported accuracies are computed on the
held-out test split.

The classifier under attack is the Adaptive Wavelet Network (AWN)~\cite{li2023awn},
a pretrained model that achieves 92.6\% classification accuracy at SNR $\geq 0$\,dB
on the clean test set.  We use the publicly available pretrained checkpoint without
any adversarial fine-tuning, so the model represents a realistic deployment target
that has not been hardened against adversarial inputs.

\subsection{Attack Configurations}

We evaluate five attacks drawn from two families: single-step gradient attacks
(FGSM), iterative gradient attacks (PGD), and optimization-based attacks
(CW-L2, EAD-L1, EAD-EN). All attacks are implemented via the torchattacks
library~\cite{kim2020torchattacks} using a \texttt{Model01Wrapper} that maps
IQ signals into the $[0,1]$ range expected by torchattacks and returns logits
only.  We use per-sample min-max normalization (\texttt{ta\_box=minmax}) to
map each sample to $[0,1]$ so that epsilon is expressed relative to the
sample's own dynamic range---a natural choice for IQ signals whose peak-to-peak
amplitude varies widely across modulations and SNR levels.

\begin{table}[t]
\caption{Attack configurations. All attacks target untargeted misclassification
  in a white-box setting. Epsilon values are expressed after per-sample min-max
  normalization to $[0,1]$.}
\label{tab:attacks}
\centering
\small
\begin{tabular}{llll}
\toprule
Attack & Norm & Key Parameters & Source \\
\midrule
FGSM   & $\ell_\infty$ & $\varepsilon = 0.03$                     & \cite{goodfellow2015fgsm} \\
PGD    & $\ell_\infty$ & $\varepsilon = 0.03$, steps $= 10$, $\alpha = 0.01$ & \cite{madry2018pgd} \\
CW     & $\ell_2$      & $c = 1.0$, steps $= 100$, lr $= 0.01$   & \cite{carlini2017towards} \\
EAD-L1 & $\ell_1$      & $c_0 = 0.01$, $\beta = 0.01$, steps $= 100$ & \cite{chen2020ead} \\
EAD-EN & elastic net   & $c_0 = 0.01$, $\beta = 0.01$, steps $= 100$ & \cite{chen2020ead} \\
\bottomrule
\end{tabular}
\end{table}

Table~\ref{tab:attacks} summarizes the attack hyperparameters.  For FGSM and
PGD, $\varepsilon = 0.03$ after min-max normalization corresponds to roughly
60\% of the average per-sample dynamic range at high SNR---a moderate threat
that leaves the signal recognizable but is effective enough to reduce accuracy
substantially.  For CW and EAD, the optimization constant $c$ controls the
trade-off between perturbation minimization and attack success; we use $c=1.0$
and $c_0=0.01$, respectively, as these yield near-zero undefended accuracy on
most modulations at high SNR while keeping perturbation norms in a physically
plausible range.

\subsection{Defense Parameters}

We evaluate nine defense configurations: two proposed adaptive spectral methods
(Adaptive-K, Spectral-Gated) and seven baselines (Kalman, Wiener,
Savitzky-Golay, Gaussian, FIR, and Randomized Smoothing), plus the undefended
baseline.  The proposed methods operate entirely in the frequency domain without
retraining.

\begin{table}[t]
\caption{Defense configurations. Classical filter parameters are calibrated
  per-SNR via grid search on the validation set; representative values at
  SNR $= 10$\,dB are shown. Proposed methods (Adaptive-K and Spectral-Gated)
  require no per-SNR calibration.}
\label{tab:defenses}
\centering
\small
\begin{tabular}{llp{4.5cm}}
\toprule
Defense & Type & Key Parameters (SNR $= 10$\,dB) \\
\midrule
Adaptive-K      & Proposed & $K$ auto per sample; energy threshold $\eta=0.95$ \\
Spectral-Gated  & Proposed & $K=50$; spectral flatness threshold $= 0.5$ \\
Kalman          & Classical & $Q=0.005$, $R=0.001$; CPU sequential \\
Wiener          & Classical & noise $= 0.01$, filter len $= 3$; freq.\ domain \\
Savitzky-Golay  & Classical & window $= 5$, order $= 1$; time domain \\
Gaussian        & Classical & $\sigma=0.5$; GPU depthwise conv1d \\
FIR             & Classical & cutoff $= 0.15$, taps $= 63$; GPU depthwise conv1d \\
Rand.\ Smoothing & Baseline  & $\sigma=0.1$, $k=1$ copies; certified robustness \\
\bottomrule
\end{tabular}
\end{table}

Table~\ref{tab:defenses} summarizes the defense parameters.  For the five
classical filters (Kalman, Wiener, Savitzky-Golay, Gaussian, FIR), we perform
a grid search over the parameter space at each SNR level using the validation
set with CW attack.  The best-performing parameter combination per SNR is stored
in a calibration file (\texttt{calibration\_params.json}) and loaded at test
time, giving the classical baselines every possible advantage.  Randomized
Smoothing~\cite{cohen2019certified} adds Gaussian noise ($\sigma=0.1$) before
classification using a single forward pass ($k=1$), consistent with its use
as an inference-time defense rather than a certified smoothed classifier.

\subsection{Evaluation Metrics}

We report the following metrics:

\textbf{Per-SNR accuracy}: fraction of test samples correctly classified at
each SNR level $s \in \{0, 2, \ldots, 18\}$ dB.  We focus on non-negative SNR
because at $\text{SNR} < 0$\,dB the clean accuracy is below 50\% and adversarial
perturbations are masked by channel noise.

\textbf{Weighted average accuracy}: the mean per-SNR accuracy, weighted by the
number of test samples at each SNR level.  Since RML2016.10a is balanced, this
reduces to the simple mean over the 10 SNR levels $\geq 0$\,dB evaluated.

\textbf{Perturbation budget curves}: accuracy as a function of attack strength
($\varepsilon$ for FGSM/PGD; optimization constant $c$ for CW/EAD) to assess
defense robustness across the threat spectrum.

\textbf{Confusion matrices}: per-class classification accuracy before and after
defense at a fixed SNR (18\,dB) to identify which modulation classes benefit
most from the proposed methods.

All evaluations use 200 samples per (SNR, modulation) cell from the test split,
yielding 2{,}200 samples per SNR level (200 samples $\times$ 11 modulation classes).

\section{Results and Analysis}
\label{sec:results}

% TODO: content in Plan 02

\subsection{Defense Comparison Overview}
% TODO: Table of all 9 defenses x 5 attacks, weighted avg accuracy
% Reference Figure \ref{fig:defense_overview}

\subsection{Accuracy vs.\ SNR for CW and EAD Attacks}
% TODO: Line plots showing per-SNR defense accuracy
% Reference Figures \ref{fig:snr_cw}, \ref{fig:snr_eadl1}

\subsection{Confusion Matrix Analysis}
% TODO: Before/after defense confusion matrices for CW and EAD at SNR=18
% Reference Figures \ref{fig:confmat_cw}, \ref{fig:confmat_eadl1}

\subsection{Perturbation Budget Analysis}
% TODO: Accuracy vs. epsilon/c curves for FGSM and CW
% Reference Figures \ref{fig:budget_fgsm}, \ref{fig:budget_cw}

\subsection{Frequency-Domain Analysis}
% TODO: FFT spectra clean vs. attacked vs. recovered
% Reference Figure \ref{fig:freq_spectra}

\subsection{Latency and Computational Cost}
% TODO: GPU latency table for all defenses; real-time feasibility claim

\section{Conclusion}
\label{sec:conclusion}

We presented an adaptive spectral defense framework for protecting
deep-learning-based automatic modulation classifiers against adversarial attacks.
Our framework requires no retraining of the classifier and incurs sub-millisecond
overhead, making it suitable for always-on deployment in real-time spectrum
monitoring and cognitive radio applications.

Through systematic evaluation on the RML2016.10a dataset against five attacks
(CW, EAD-L1, EAD-EN, FGSM, PGD) and comparison with five calibrated classical
signal processing baselines, we established four principal findings:

\begin{enumerate}

\item \textbf{Adversarial attacks on AMC are control-plane attacks.}  CRC
  integrity analysis shows that adversarial perturbations collapse AMC accuracy
  to near zero while leaving the underlying signal data largely intact (CRC pass
  rate $> 71\%$ with oracle demodulation).  This characterization motivates
  recovery as the correct defense strategy: the perturbation need only be
  removed to restore correct classification.

\item \textbf{Adaptive-K FFT filtering outperforms all classical baselines on
  optimization-based attacks.}  By automatically selecting the number of
  retained spectral components per sample based on cumulative energy threshold
  ($\eta = 0.95$), Adaptive-K achieves 77.4\% weighted-average accuracy against
  CW (vs.\ 72.7\% for the best classical filter, Kalman) and 79.8\% against
  EAD-L1 (vs.\ 72.5\% for Kalman)---improvements of 4.7 and 7.3 percentage
  points, respectively.

\item \textbf{Classical signal processing filters are ineffective against
  optimization-based attacks despite per-SNR calibration.}  Savitzky-Golay,
  Wiener, and Kalman filters all perform at or below the undefended baseline on
  CW and EAD attacks, because their smoothing operations spread spectrally
  concentrated adversarial energy rather than removing it.  This negative result
  is practically important: it demonstrates that deploying conventional filtering
  as a security measure provides a false sense of protection against
  optimization-based adversaries.

\item \textbf{Spectral-gated routing improves recovery for wideband analog
  modulations.}  By bypassing Top-K filtering for samples whose spectral
  flatness exceeds a threshold (indicating wideband energy distribution),
  Spectral-Gated avoids distorting AM-SSB and WBFM signals where Top-K alone
  introduces classification-relevant artifacts.  Spectral-Gated achieves
  76.1\% on CW vs.\ 77.4\% for Adaptive-K, offering a useful accuracy--overhead
  trade-off for heterogeneous signal environments.

\end{enumerate}

Several directions remain open for future work.  First, the current evaluation
assumes a defense-unaware attacker; an adversary with knowledge of the Top-K
operator could craft perturbations that survive the frequency mask.  Evaluating
such \emph{adaptive attacks}---where the attack optimization incorporates the
defense---is an important robustness validation that we leave to future work.
Second, we focused on RML2016.10a (128-sample signals, 11 classes); extending
the evaluation to RML2018.01a (1024-sample signals, 24 classes) would establish
the generality of the spectral defense approach across longer signals and a
richer modulation set.  Third, combining adaptive spectral filtering with
adversarial training could provide complementary robustness: training-time
regularization hardens the classifier's decision boundaries while
inference-time filtering removes perturbation energy before the signal reaches
the network.


% ---- Bibliography ----
\bibliographystyle{IEEEtran}
\bibliography{refs}

\end{document}
